\chapter{The source cone: two closed faces and two contact families}
\label{ch:geometry}
\label{component:p0c8f420148b7}
\label{p0c8f420148b7:sec:physical}
\label{p0c8f420148b7:eq:parent}
\label{p0c8f420148b7:prop:compat}
\label{component:pd9a6195f0af3}
\label{pd9a6195f0af3:sec:closed-faces-short}
\label{pd9a6195f0af3:lem:closed-face-orbits}
\label{pd9a6195f0af3:lem:quadratic-matrix-lift}
\label{pd9a6195f0af3:thm:closed-face-adjacent}
\label{pd9a6195f0af3:thm:closed-face-opposite}
\label{pd9a6195f0af3:cor:all-closed-faces}
\section{The exact physical polynomial}

Use the affine functional $H=1-(E_{00}-E_{01}+E_{10}+E_{11})/2$.
Order the eight losing events as
\[
 (\lambda,n,m,C,A,D,\alpha,\beta)
 =\bigl((++|01),(+-|00),(+-|11),(-+|10),
             (-+|00),(-+|11),(--|01),(+-|10)\bigr).
\]
Every event vanishes at the chosen PR point.  Each is one at its own
deterministic CHSH-facet neighbour and zero at the other seven.
Consequently every affine functional has the unique expansion
\[
 f=dH+4\sum_i a_i p_i,\qquad
 d=-f(\mathrm{PR}),\qquad a_i=f(v_i)/4.                 \tag{SG1}
\]
For a quantum-valid functional negative at that PR point, $d>0$ and
$a_i\geq0$.  If a quantum-valid functional is nonnegative at every PR
point, it is nonnegative on the entire nonsignalling polytope and is a
nonnegative combination of its probability inequalities.  The latter
sources have immediate separated square certificates.  Here is a complete short vertex argument. The nonsignalling affine
space has dimension eight. If no local marginal is deterministic, each
two-by-two outcome table has positive row and column sums, hence at
most two zero entries. A vertex requires at least eight active
probability inequalities. Thus every table has exactly two zeros,
necessarily in opposite corners. Write its perfect correlation or
anticorrelation as $a_x=s_{xy}b_y$, with $s_{xy}=\pm1$.
If $\prod_{x,y}s_{xy}=-1$, these equations force all four marginals
to vanish and give a PR vertex. If the product is $+1$, one marginal
remains free and the table is a mixture of two deterministic points,
so an interior marginal value is not a vertex. If instead one
marginal is deterministic, say $A_0$, a local model is obtained by
sampling $A_1$ and then sampling $B_0,B_1$ independently conditional
on $A_1$, using the two given joint tables. Nonsignalling makes their
$A_1$ marginals agree; the tables involving fixed $A_0$ are reproduced
automatically. Such a vertex is therefore deterministic. This proves
that the only vertices are the sixteen deterministic points and the
eight PR points. Linear separation (equivalently finite-dimensional
Farkas duality) now gives the asserted event representation.

Put $D_u=1+u^2$, $D_v=1+v^2$, and write the real state as
$\psi=(W,X,Y,Z)$.  The entire cleared parent is
\[
\begin{split}
 F={}&4\lambda D_u(W+vX)^2+4nD_uD_vX^2
       +4m(vW-X+uvY-uZ)^2\\
 &+4CD_v(uW-Y)^2+4AD_uD_vY^2
 +(4D+2d)(uW+uvX-Y-vZ)^2\\
 &+(4\alpha+2d)D_u(vY-Z)^2
 +(4\beta+2d)D_v(X+uZ)^2-2dD_uD_vZ^2.
\end{split}                                                   \tag{SG2}
\]
It is $D_uD_v\|\psi\|^2f(q)$ for a nonzero unnormalised state.
All these factors are pure testing factors.  Write $P(u,v)$ for its
symmetric state matrix.  If $d>0$ and $uv\ne0$, its $Z=0$ principal
block is positive definite, since the signed part alone is
\[
 2d\{(uW+uvX-Y)^2+D_uv^2Y^2+D_vX^2\}.                 \tag{SG3}
\]
Thus $P$ has at least three positive eigenvalues, its kernel at a
zero is one-dimensional with $Z\ne0$, and
\[
                 P\succeq0\quad\Longleftrightarrow\quad\det P\geq0.
                                                               \tag{SG4}
\]

Here is the entire determinant datum needed for the geometric
reduction.  Set
\[
\begin{array}{ll}
t_1=d+2(\lambda+n+D),&t_2=d+2(\alpha+m+A),\\
t_3=d+2(\lambda+C+D),&t_4=d+2(\alpha+\beta+m),\\
t_5=d+2(\lambda+\beta+n),&t_6=d+2(\alpha+C+A),\\
t_7=d+2(\beta+n+m),&t_8=d+2(C+A+D).
\end{array}
\]
All $t_i>0$.  Directly expanding the four-dimensional determinant of
SG2, or applying Cauchy--Binet to its eight positive and one negative
rank-one packets, gives
\[
 \det P=16D_u^2D_v^2u^2v^2L(u^2,v^2),\qquad
 L(U,V)=h+\sum_{(i,j)\ne(1,1)}q_{ij}U^{i-1}V^{j-1},       \tag{SG5}
\]
where
\[
\begin{array}{ll}
q_{00}=4\lambda\alpha t_7t_8,&
q_{01}=4(\lambda Dt_7t_6+\alpha mt_5t_8),\\
q_{02}=4mDt_5t_6,&
q_{10}=4(\alpha At_7t_3+\lambda nt_8t_4),\\
q_{12}=4(\beta mt_1t_6+CDt_5t_2),&
q_{20}=4nAt_3t_4,\\
q_{21}=4(\beta nt_3t_2+CAt_1t_4),&
q_{22}=4\beta Ct_1t_2.
\end{array}                                                   \tag{SG6}
\]
The central coefficient is fixed without a long expansion by
$h=\det P(1,1)/256-\sum_{(i,j)\ne(1,1)}q_{ij}$.
The eight nonnegative outer coefficients determine the possible
contact geometries.

\section{The twelve elementary angular identities}
\label{sec:angular-circuits}

The determinant of the state matrix has eight nonconstant Laurent
monomials. Their coefficients are nonnegative. A stationary zero
therefore has a particularly simple algebraic decomposition, which
also explains the cubic and quartic identities used below.

Let $E=\{-1,0,1\}^2\setminus\{(0,0)\}$ and write
\[
 L(U,V)=h+\sum_{e\in E}q_eU^{e_1}V^{e_2},\qquad q_e\ge0.
\]
Suppose $U_c,V_c>0$ satisfy $L=L_U=L_V=0$. Put
$w_e=q_eU_c^{e_1}V_c^{e_2}$. The value and derivative equations are
exactly
\[
 h=-\sum_e w_e,\qquad \sum_e w_e e=0.
 \label{eq:angular-balance}
\]
Thus the positive coefficients balance around the origin in the
exponent plane.

\begin{lemma}[Balanced Laurent polynomial]
Let $K_c,K_t$ be labelled real closed extensions of $\mathbb R$.
If the coefficients and contact above belong to $K_c$, then for every
$U,V>0$ in $K_t$ the separated evaluation of $L$ is a finite sum of
squares in $K_c\otimes_{\mathbb R}K_t$.
\end{lemma}
\begin{proof}
First decompose the coefficient vector $w$ using its balance equation.
If $w\ne0$, the origin lies in the convex hull of its active exponents.
A minimal subset containing the origin has two or three elements.
For the eight exponents in $E$, the possibilities, up to rotation
and reflection of the square, are
\[
\begin{array}{c|c|c}
\text{exponents}&\text{positive integer weights}&\text{number}\\ \hline
\{e,-e\}&(1,1)&4\\
\{(1,1),(-1,0),(0,-1)\}&(1,1,1)&4\\
\{(1,1),(-1,1),(0,-1)\}&(1,1,2)&4.
\end{array}
\]
Indeed a two-element set must be opposite. A triangle containing the
origin contains either one corner and its two opposite axial points,
or two adjacent corners and their opposite axial point. These are
the last two rows. For its integer weight vector $b$, subtract
$\lambda b$, where $\lambda=\min_{b_e>0}w_e/b_e$.
The remainder is nonnegative and balanced, and has fewer active
entries. Repeating at most eight times gives
$w=\sum_j\lambda_jb^{(j)}$ with $\lambda_j>0$ in $K_c$.

Set $X=U_c^{-1}\otimes U$ and $Y=V_c^{-1}\otimes V$.
Their rational powers mean separate positive roots in the two fields.
For each row of the table the monomials $p_i=X^{e_{i1}}Y^{e_{i2}}$
satisfy $\prod_i p_i^{b_i}=1$. The required elementary identities are
\begin{align*}
 p+p^{-1}-2&=(\sqrt p-p^{-1/2})^2,\\
 a^3+b^3+c^3-3abc
   &=\tfrac12(a+b+c)\bigl[(a-b)^2+(b-c)^2+(c-a)^2\bigr],\\
 a^4+b^4+c^4+d^4-4abcd
   &=(a^2-b^2)^2+(c^2-d^2)^2+2(ab-cd)^2.
\end{align*}
For weights $(1,1,1)$ take $a=p_1^{1/3}$, $b=p_2^{1/3}$,
$c=p_3^{1/3}$; then $abc=1$. Each of $a,b,c$ is an elementary
tensor square, so the middle line is a sum of tensor squares.
For weights $(1,1,2)$ take $a=p_1^{1/4}$, $b=p_2^{1/4}$ and
$c=d=p_3^{1/4}$; then $abcd=1$.
Finally multiply the squares by $\sqrt{\lambda_j}\otimes1$ and sum.
Equation~\eqref{eq:angular-balance} makes the result exactly $L$.
\end{proof}

In a single real closed field this proves that a stationary zero of
$L$ is a global minimum, using only the three displayed identities.
For the Bell state matrix, its positive three-dimensional principal
block then turns determinant nonnegativity into matrix positivity.
The same calculation proves a separated sum-of-squares identity for
the determinant at independent tests. To obtain a certificate for the
whole state polynomial we must also control the state-normal terms;
the following matrix lemmas perform that step.


\section{Stationary zeros and physicality}

\paragraph{Contact lemma.}
Suppose $d>0$, all event weights are nonnegative, and there exist
positive $u_c,v_c$ and a nonzero real state $\psi_c$ such that
\[
 P(u_c,v_c)\psi_c=0,
 \qquad\psi_c^TP_u(u_c,v_c)\psi_c=0,
 \qquad\psi_c^TP_v(u_c,v_c)\psi_c=0.                    \tag{SG7}
\]
Then the source is physical at every angle and every state.  This
statement holds over every real closed field.

\emph{Proof.}
SG3 gives three positive complementary state eigenvalues. The adjugate
at contact is a positive multiple of $\psi_c\psi_c^T$, so SG7 and SG5
give $L=L_U=L_V=0$ there. The balanced Laurent identity just proved
makes $L$ nonnegative at every positive pair in the same field.
Equation SG4 therefore gives physicality for $uv\ne0$.
At a commuting or projective endpoint, simultaneous diagonalization
of the commuting local pair reduces the CHSH operator to $\pm2$
times a remote reflection, so $H\succeq0$. The event terms are
nonnegative there as well. Homogenizing SG2 includes all endpoints.
This proves the claim over every real closed field.\hfill$\square$

This lemma is also the justification for deleting a positive event
whose amplitude vanishes at a contact.  Its square and all first
derivatives vanish; deleting it preserves SG7.  It is not legitimate
to delete an event whose amplitude is nonzero.

\section{The geometry of a stationary zero}

\paragraph{Coercivity at every physical finite contact.}
Every physical source with $d>0$ and a finite noncommuting zero has
coercive Newton support.  Indeed, put
$a_e=q_eU_c^{e_1}V_c^{e_2}>0$ for the active exponent vectors.
The two contact derivatives give
\[
                         \sum_e a_e e=0.                \tag{SG17}
\]
If zero is not in the interior of their convex hull, a separating
linear functional nonnegative on every active vector must vanish
on all of them, by SG17.  Thus the active support is empty or lies
on one of the four lines through the origin in the square grid.
A nonempty one-sided subset cannot satisfy SG17.  The horizontal
opposite pair $q_{01},q_{21}>0$ has four possible product pairings
in SG6.  These respectively force $\lambda n$, $CD$, $\beta m$
or $\alpha A$ positive, activating an off-line coefficient
$q_{10}$ or $q_{12}$.  The vertical pair similarly forces
$\alpha m$, $CA$, $\beta n$ or $\lambda D$ and activates
$q_{01}$ or $q_{21}$.

For the diagonal pair $q_{00},q_{22}>0$, all off-line coefficients
can vanish only on the alternating matching
$\lambda,\alpha,\beta,C>0$, $n=m=A=D=0$.
For the other diagonal they can vanish only on the opposite matching
$n,m,A,D>0$, $\lambda=C=\alpha=\beta=0$.
The latter is not physical: SG2 at $u=v=1$ and
$\psi=(1,-\varepsilon,\varepsilon,1)$ gives
\[
 F=16\varepsilon\{(d+n+m+A+D)\varepsilon-d\}<0
 \quad\text{for }\quad
 \varepsilon=\frac{d}{2(d+n+m+A+D)}.
\]
The first matching is its physical one-step cycle rotation.
If every outer coefficient is zero, SG6 forbids every adjacent or
opposite positive pair in the cycle
$(\lambda,D,C,A,\alpha,m,\beta,n)$.
Any maximal remaining set is a triple with cyclic gaps $(2,3,3)$;
smaller sets extend to such a triple.  The eight triples are rotations
of $\{\lambda,m,A\}$.  At $u=v=1$, $\psi=(1,-1,0,2)$
annihilates precisely the needed three event rows and has signed-core
value $-4$, giving $F=-4d<0$.  This excludes empty support too.

For completeness, the physical cycle symmetries used here are the
CHSH-preserving maps on event coordinates $(x,y,a,b)$
\[
\begin{split}
g_A:(x,y,a,b)&\mapsto(1-x,y,a,(-1)^yb),\\
g_B:(x,y,a,b)&\mapsto(x,1-y,(-1)^{1-x}a,b),\\
g_T:(x,y,a,b)&\mapsto(y,x,(-1)^yb,(-1)^xa).
\end{split}
\]
In the displayed cycle, indexed from zero, they act as
$i\mapsto5-i,3-i,2-i$ modulo eight.  Their compositions give
a one-step rotation and a reflection.  The assertions follow
by inserting the eight event definitions in SG1.

\paragraph{Strict single-field contact Hessian.}
The logarithmic angular Hessian is now
\[
                 \mathcal H_L=\sum_e a_e ee^T\succ0,       \tag{SG18}
\]
since the active vectors span the plane.  Write the state matrix
as $P=\left(\begin{smallmatrix}H_3&b\\b^T&c\end{smallmatrix}\right)$
in the coordinates $(W,X,Y;Z)$ and set $Z=1$.
The positive definite state block SG3 gives, within the single
source field,
\[
 F=(z+H_3^{-1}b)^TH_3(z+H_3^{-1}b)+s(u,v),
 \qquad s=\det P/\det H_3,\quad z=(W,X,Y)^T.             \tag{SG19}
\]
At contact $s=0$ and its gradient vanishes.  Its angular Hessian
is the positive scalar $16D_u^2D_v^2u^2v^2/\det H_3$ times
the congruence of SG18 by $\operatorname{diag}(2/u,2/v)$.
Consequently the full Hessian of $F$ in $(W,X,Y,u,v)$ has
positive state block $2H_3$ and a positive angular Schur complement;
it is positive definite.  All formulas for derivatives and
positive definiteness are algebraic.  The finite real assertion
therefore transfers to every real closed field; logarithms are
only notation for $U\partial_U,V\partial_V$.
The inverse in SG19 belongs to the source field.

\paragraph{Positive contact state orientation.}
Let $F_H$ denote SG2 with $d=1$ and all eight event weights zero.
Its state matrix satisfies
\[
              \det P_H=-16D_u^2D_v^2u^2v^2\ne0.          \tag{SG20}
\]
This follows by a four-dimensional determinant expansion of the
three signed-core squares and the negative $Z^2$ term in SG2.
If all positive-weight event rows vanished at a physical contact,
their state matrices would annihilate the contact state.  The
kernel equation would then give $P_H\psi_c=0$, contradicting
SG20.  Hence at least one event contributes positively.  Since
$F(\psi_c)=0$ and $d>0$, the core value $F_H(\psi_c)$ is strictly
negative.

For positive $u,v$, the $X,Y$ block of $F_H$ is
$2\left(\begin{smallmatrix}h_0&-uv\\-uv&h_0\end{smallmatrix}\right)$,
where $h_0=1+v^2+u^2v^2$.  Its Schur form is exactly SG15.
It is nonnegative when $WZ\le0$.  Thus the negative core value
forces
\[
                  W_cZ_c>0,\qquad Z_c=1\Longrightarrow W_c>0.
                                                               \tag{SG21}
\]

\paragraph{The forbidden-deletion implication.}
If a positive event row vanished at a contact, deleting that
coefficient would preserve the four kernel equations and both
angular derivatives.  SG7 would make the resulting nonnegative
source physical, with the same finite contact.  SG17--SG18
therefore make its support coercive.  We obtain the exact rule
\[
 \text{a positive contact event cannot vanish if deleting it
 makes the Newton support noncoercive}.                       \tag{SG22}
\]
In particular every deletion introducing an exclusive omitted pair
is forbidden by SG9.  This rule will fix the signs of the contact coordinates in
Chapter~\ref{ch:entry}.

\section{Reduction to two contact families}

The active exponent vectors in SG5 lie on the boundary of
$\{-1,0,1\}^2$.  The exponential sum is coercive exactly when zero
is in the interior of their convex hull.  One direction follows from
uniform positive support in every unit direction; the other follows
from a separating direction along which every active exponential is
nonincreasing.

For this eight-point configuration a separating direction can be
chosen from $(\pm1,0),(0,\pm1),(\pm1,\pm1)$.  For example the axis
condition in SG6 gives
$\lambda\alpha=\lambda D=\alpha m=mD=0$, whose minimal omitted
event pairs are $\{\lambda,m\}$ and $\{\alpha,D\}$.
For a diagonal the additional minimal pair is
$\{\lambda,\alpha\}$; the remaining minimal triples already
contain an axis pair.  Rotating the event cycle
\[
                 (\lambda,D,C,A,\alpha,m,\beta,n)          \tag{SG8}
\]
gives exactly the twelve pairs of exclusive events.  Therefore
\[
 \text{coercive exponent support}\quad\Longleftrightarrow\quad
 \text{every two omitted losing events are compatible}.       \tag{SG9}
\]
Compatibility means agreement on each shared local measurement.
Pairwise compatible events extend to a single deterministic
assignment.  Such an assignment loses one or three CHSH inputs;
there are consequently at most three compatible omissions.
In SG8 compatible pairs have cyclic distance one or two, and a
compatible triple is three consecutive vertices.

\paragraph{Extreme-source lemma.}
Normalise the canonical source by $d+\sum a_i=1$ and impose
physicality and nonnegativity.  This is a compact convex subset of
an eight-dimensional simplex.  Every extreme point with $d>0$ and
coercive support has a finite stationary zero and exactly five or
six positive event weights.

\emph{Proof.}
Coercivity gives a finite minimum $\mu\geq0$ of $L$.  If $\mu>0$,
take sufficiently small opposite nonradial perturbations inside the
exact active coefficient face, preserving the normalisation.  Every
positive $q_e$ changes by relative error at most $\varepsilon$,
every zero $q_e$ remains zero, and $h$ changes by absolute error at
most $\varepsilon$.  Thus, uniformly at all positive angles,
\[
 L'\geq(1-\varepsilon)L+\varepsilon h-\varepsilon
       \geq(1-\varepsilon)\mu+\varepsilon h-\varepsilon>0.
\]
The perturbations remain physical by SG4 and the direct endpoint
argument, contradicting extremality.  The perturbation exists:
SG9 permits at most three omissions, hence at least six active
coordinates including $d$.  Thus $\mu=0$ and there is a contact.

Fix its angles and state.  SG7 consists of six homogeneous linear
equations in the source coefficients.  If at least seven events
were positive, there would be at least eight active coordinates;
their solution space has dimension at least two and contains the
source.  Intersecting its directions with coefficient sum zero
leaves a nonzero direction.  Small opposite perturbations preserve
all signs and SG7, and are physical by the contact lemma.  This again
contradicts extremality.  Hence at most six events are positive.
The lower bound five follows from SG9. \hfill$\square$

The CHSH-preserving event relabellings generate the dihedral group
of SG8.  Thus the coercive extreme sources, including every allowed
third omission, reduce to the two patterns
\[
\begin{array}{ll}
\mathsf C_1:&\lambda=n=0,\quad d,m,C,A,\alpha>0,
                  \quad D,\beta\geq0,\quad D+\beta>0,\\
\mathsf C_2:&\lambda=\beta=0,\quad d,m,C,A,D,\alpha>0,
                  \quad n\geq0.
\end{array}                                                     \tag{SG10}
\]
The C1 third omission is $D=0$ or $\beta=0$; the C2 third omission
is $n=0$, already a C1 orbit.  This overlap is harmless.  It is not
legitimate to replace positive infinitesimal coefficients by zeros.

Every point of the normalised source body is a convex combination
of at most nine extreme points.  The fixed nine-term assertion,
the midpoint definition of extremality, SG7, and physicality all
have first-order polynomial formulations.  Their real proof
therefore transfers to a real closed coefficient field, without
using topological compactness or arbitrary infinite decompositions
in that field.  A nonnegative decomposition preserves each literal
zero coefficient term by term.

\section{Noncoercive sources have explicit lifts}

Every noncoercive source omits an exclusive pair.  Up to the same
finite symmetries there are only two possibilities:
$\lambda=m=0$ and $\lambda=\alpha=0$.

We first record the elementary quadratic matrix lemma.  A real
symmetric matrix polynomial $K(t)=K_0+tK_1+t^2K_2$ is PSD for every
real $t$ exactly when there is a skew-symmetric $S$ such that
\[
 \begin{pmatrix}K_0&K_1/2+S\\K_1/2-S&K_2\end{pmatrix}\succeq0.
                                                               \tag{SG11}
\]
For completeness, the cone of Gram-representable coefficient
triples is closed because the trace of a representing PSD Gram is
$\operatorname{tr}(K_0+K_2)$.  A nonnegative dual functional has
block Hankel matrix $\left(\begin{smallmatrix}U&V\\V&W\end{smallmatrix}\right)
\succeq0$, with $U,V,W$ symmetric.  Diagonalising
$U^{+/2}VU^{+/2}$ on $\operatorname{ran}U$ writes its moment part as
$\sum_j(z_j,t_jz_j)(z_j,t_jz_j)^T$; the remainder is
$\operatorname{diag}(0,W-VU^+V)\succeq0$.
Its value on $K$ is consequently
$\sum_jz_j^TK(t_j)z_j+\sum_iw_i^TK_2w_i\geq0$.
Separation proves SG11, including singular $U$.  The fixed-size
statement transfers to real closed fields.

For $\lambda=m=0$, use SG11 with
\[
\begin{split}
 K_0&=\operatorname{diag}(4C+4D+2d,4A,4\alpha),\\
 K_2&=\operatorname{diag}(4C,4A+4\alpha+2d,4D),\\
 K_1&=\begin{pmatrix}0&0&-(4D+2d)\\
                    0&0&-(4\alpha+2d)\\
                    -(4D+2d)&-(4\alpha+2d)&0\end{pmatrix}.
\end{split}                                                     \tag{SG12}
\]
Necessity follows by setting $W=a/u$ and letting $u\to0$ in the
homogeneous state form.  The surviving form, after $R=a-Y$, is
$2t_5D_vX^2+(R,Y,Z)^TK(v)(R,Y,Z)$.
For sufficiency the determinant identity gives
\[
 \det K(v)=\frac{8D_v}{t_5}(q_{10}+hv^2+q_{12}v^4),
\]
whereas
$L=h+q_{10}/V+q_{12}V+U(q_{20}/V+q_{21}+q_{22}V)$.
PSD of $K$ makes the first three terms nonnegative; all remaining
terms are nonnegative.  SG4 and the endpoint argument finish the
equivalence.  When $d=0$ all forms are event-square sums.
The Gram is strictly feasible at
$C=A=2$, $D=\alpha=d=n=\beta=1$, $S=0$:
after permutation its two Schur margins are $125/56>0$.

For $\lambda=\alpha=0$, put $p=4C+4D+2d$,
$q=4m+4\beta+2d$, $r=4D+2d$.  The exact lift is
\[
 B_1=\begin{pmatrix}4m&-4m&\xi\\-4m&q+4n&q\\\xi&q&q\end{pmatrix}
 \succeq0,\quad
 B_2=\begin{pmatrix}p+4A&4D&-p\\
                 4D&4D&-(4m+r)-\xi\\
                 -p&-(4m+r)-\xi&p\end{pmatrix}\succeq0.          \tag{SG13}
\]
Indeed set $u=\varepsilon s$, $v=\varepsilon$,
$(W,Z)=(w,z)/\varepsilon$, $X=x$, $Y=y$.
The surviving form is
\[
\begin{split}
 G_s={}&4nx^2+4m(w-x-sz)^2+4C(sw-y)^2+4Ay^2\\
       &+r(sw-y-z)^2+(4\beta+2d)(x+sz)^2-4dyz-2dz^2.
\end{split}
\]
Its determinant is $16(q_{02}+hs^2+q_{20}s^4)$.  The same
nonnegative-remainder argument proves equivalence with physicality.
Applying SG11 to its state matrix, the zero diagonal entries on
$sx,sy$ force their Gram rows to vanish.  Coefficient matching leaves
only $\xi$, and the remaining Gram splits into the two blocks SG13.
Equivalently
$G_s=(w,x,sz)B_1(w,x,sz)^T+(y,z,sw)B_2(y,z,sw)^T$;
the $\xi$ terms cancel because $w(sz)=z(sw)$.
At $C=A=n=\beta=2$, $m=D=d=1$, $\xi=-5$, the successive
principal minors are $(4,72,234)$ and $(22,72,234)$.
Both lifts retain all nonnegative coefficient constraints.

\paragraph{Why these source lifts give separated certificates.}
Represent either source cone as $a\geq0$ and
$\exists z\,M(a,z)\succeq0$, absorbing the signs into diagonal
blocks.  The displayed points make this entire homogeneous lift
strictly feasible.  The image of its dual PSD cone is closed:
pairing with a fixed strict feasible matrix bounds the trace of
every PSD preimage of a convergent dual sequence.  Separation
therefore proves attained dual representation for every functional
nonnegative on the source cone.  This is a finite-dimensional
argument, and its fixed matrix-size quantified conclusion transfers.

Given a source in $K_c$, choose its lift $A_c\succeq0$ there.
Given a physical test in $K_t$, its evaluation functional belongs
to the dual source cone in $K_t$; choose its attained dual matrix
$B_t\succeq0$ there.  Their separated trace pairing equals the
slack because all auxiliary-coordinate coefficients cancel.
Factoring $A_c=RR^T$ and $B_t=SS^T$ in their respective fields gives
\[
 \langle A_c,B_t\rangle_{\otimes}
       =\sum_{i,j}\left(\sum_k R_{ki}\otimes S_{kj}\right)^2.
                                                               \tag{SG14}
\]
No mixed inverse or limiting certificate is involved.  Literal
zero coefficients and singular matrices are included.

\section{The remaining quadratic form}

Let $K_c,K_t$ be separately labelled real closed extensions of
$\mathbb R$ and $T=K_c\otimes_{\mathbb R}K_t$.
The source partition in Chapter~\ref{ch:cover} reduces the following
assertion to explicit local calculations.

\paragraph{Canonical assertion $\mathcal C$.}
Take $d=1$ and source coefficients in either pattern SG10.
Assume only that there exist $u_c,v_c>0$, $W_c>0$, and
$X_c,Y_c\in K_c$ for which SG7 holds with
$\psi_c=(W_c,X_c,Y_c,1)$.  For every independent
$u,v,W>0$ and $X,Y\in K_t$, the full expression SG2 with
$Z=1$ belongs to $\Sigma T^2$.

Physicality follows from SG7. Every source produced by the extreme
reduction has this form after the angle signs and state normalization
are fixed in its own field. The independent testing state remains
arbitrary in $\mathcal C$.

\paragraph{Assembly lemma.}
Once $\mathcal C$ is proved, every transferred quantum-valid affine
functional and every independent transferred quantum behaviour
have separated slack in $\Sigma T^2$.  Therefore the universal
tensor-SOS finite-lift criterion applies to the full eight-coordinate
quantum set.

\emph{Proof.}
The source reduction gives at most nine canonical source atoms,
each of which is an event atom, a closed-face atom, or SG10 with
an actual contact.  Event atoms and closed faces were handled above.
Restore $d$ by multiplying each root by $\sqrt d\otimes1$.

For a canonical contact source, finite angle signs are harmless.
The polynomial is unchanged by
\[
 (u;W,X,Y,Z)\mapsto(-u;W,X,-Y,-Z),\qquad
 (v;W,X,Y,Z)\mapsto(-v;W,-X,Y,-Z).
\]
Literal zero/infinite angles have nonnegative pure-testing $H$ and
nonnegative event squares, as already proved.  For $u,v>0$, the
signed core is nonnegative whenever $WZ\leq0$.  To check this,
put $h_0=1+v^2+u^2v^2$; its $X,Y$ block is
$2\left(\begin{smallmatrix}h_0&-uv\\-uv&h_0\end{smallmatrix}\right)
\succ0$, and the remaining Schur form on $W,Z$ is
\[
 \frac{2D_uD_v}{h_0^2-u^2v^2}
 \{u^2v^2(W^2+Z^2)-2uvh_0WZ\}\geq0.                 \tag{SG15}
\]
This is a pure testing calculation.  Its square root in $K_t$,
together with $\sqrt d$ in $K_c$, certifies the signed core;
every event remains a square.  It includes $Z=0$.
Otherwise $Z\ne0$ and $W/Z>0$; divide the state by $Z$ inside
$K_t$ and apply $\mathcal C$.  Restore $Z^2$ afterwards.
Real and imaginary state parts add because the measurements are
real.  All homogeneous angle normalisations are pure testing
squares and have nonzero denominators on their stated charts.

Finally every testing behaviour is a mixture of at most nine real
two-qubit behaviours.  If source and testing weights are $a_j,w_i$
and atom certificates are $\sum_kh_{ijk}^2$, the entire slack is
\[
       \sum_{i,j,k}
           [ (\sqrt{a_j}\otimes\sqrt{w_i})h_{ijk}]^2.
\]
The affine constant is preserved because $\sum_iw_i=1$.
Zero weights are never inverted.  This proves the universal tensor
assertion; applying the criterion in Chapter~\ref{ch:tensor} is the final
qualitative existence step. \hfill$\square$

